
\documentclass[a4paper,11pt]{article}
\usepackage[a4paper, margin=8em]{geometry}

% usa i pacchetti per la scrittura in italiano
\usepackage[french,italian]{babel}
\usepackage[T1]{fontenc}
\usepackage[utf8]{inputenc}
\frenchspacing 

% usa i pacchetti per la formattazione matematica
\usepackage{amsmath, amssymb, amsthm, amsfonts}

% usa altri pacchetti
\usepackage{gensymb}
\usepackage{hyperref}
\usepackage{standalone}

% imposta il titolo
\title{Appunti Comunicazioni Numeriche}
\author{Luca Seggiani}
\date{2026}

% imposta lo stile
% usa helvetica
\usepackage[scaled]{helvet}
% usa palatino
\usepackage{palatino}
% usa un font monospazio guardabile
\usepackage{lmodern}

\renewcommand{\rmdefault}{ppl}
\renewcommand{\sfdefault}{phv}
\renewcommand{\ttdefault}{lmtt}

% disponi il titolo
\makeatletter
\renewcommand{\maketitle} {
	\begin{center} 
		\begin{minipage}[t]{.8\textwidth}
			\textsf{\huge\bfseries \@title} 
		\end{minipage}%
		\begin{minipage}[t]{.2\textwidth}
			\raggedleft \vspace{-1.65em}
			\textsf{\small \@author} \vfill
			\textsf{\small \@date}
		\end{minipage}
		\par
	\end{center}

	\thispagestyle{empty}
	\pagestyle{fancy}
}
\makeatother

% disponi teoremi
\usepackage{tcolorbox}
\newtcolorbox[auto counter, number within=section]{theorem}[2][]{%
	colback=blue!10, 
	colframe=blue!40!black, 
	sharp corners=northwest,
	fonttitle=\sffamily\bfseries, 
	title=~\thetcbcounter: #2, 
	#1
}

% disponi definizioni
\newtcolorbox[auto counter, number within=section]{definition}[2][]{%
	colback=red!10,
	colframe=red!40!black,
	sharp corners=northwest,
	fonttitle=\sffamily\bfseries,
	title=~\thetcbcounter: #2,
	#1
}

% U.D.M
\newcommand{\amp}{\ensuremath{\mathrm{A}}}
\newcommand{\volt}{\ensuremath{\mathrm{V}}}
\newcommand{\meter}{\ensuremath{\mathrm{m}}}
\newcommand{\second}{\ensuremath{\mathrm{s}}}
\newcommand{\farad}{\ensuremath{\mathrm{F}}}
\newcommand{\henry}{\ensuremath{\mathrm{H}}}
\newcommand{\siemens}{\ensuremath{\mathrm{S}}}

% circuiti
\usepackage{circuitikz}
\usetikzlibrary{babel}

% disegni
\usepackage{pgfplots}
\pgfplotsset{width=10cm,compat=1.9}

% disponi codice
\usepackage{listings}
\usepackage[table]{xcolor}

\lstdefinestyle{codestyle}{
		backgroundcolor=\color{black!5}, 
		commentstyle=\color{codegreen},
		keywordstyle=\bfseries\color{magenta},
		numberstyle=\sffamily\tiny\color{black!60},
		stringstyle=\color{green!50!black},
		basicstyle=\ttfamily\footnotesize,
		breakatwhitespace=false,         
		breaklines=true,                 
		captionpos=b,                    
		keepspaces=true,                 
		numbers=left,                    
		numbersep=5pt,                  
		showspaces=false,                
		showstringspaces=false,
		showtabs=false,                  
		tabsize=2
}

\lstdefinestyle{shellstyle}{
		backgroundcolor=\color{black!5}, 
		basicstyle=\ttfamily\footnotesize\color{black}, 
		commentstyle=\color{black}, 
		keywordstyle=\color{black},
		numberstyle=\color{black!5},
		stringstyle=\color{black}, 
		showspaces=false,
		showstringspaces=false, 
		showtabs=false, 
		tabsize=2, 
		numbers=none, 
		breaklines=true
}

\lstdefinelanguage{javascript}{
	keywords={typeof, new, true, false, catch, function, return, null, catch, switch, var, if, in, while, do, else, case, break},
	keywordstyle=\color{blue}\bfseries,
	ndkeywords={class, export, boolean, throw, implements, import, this},
	ndkeywordstyle=\color{darkgray}\bfseries,
	identifierstyle=\color{black},
	sensitive=false,
	comment=[l]{//},
	morecomment=[s]{/*}{*/},
	commentstyle=\color{purple}\ttfamily,
	stringstyle=\color{red}\ttfamily,
	morestring=[b]',
	morestring=[b]"
}

% disponi sezioni
\usepackage{titlesec}

\titleformat{\section}
	{\sffamily\Large\bfseries} 
	{\thesection}{1em}{} 
\titleformat{\subsection}
	{\sffamily\large\bfseries}   
	{\thesubsection}{1em}{} 
\titleformat{\subsubsection}
	{\sffamily\normalsize\bfseries} 
	{\thesubsubsection}{1em}{}

% disponi alberi
\usepackage{forest}

\forestset{
	rectstyle/.style={
		for tree={rectangle,draw,font=\large\sffamily}
	},
	roundstyle/.style={
		for tree={circle,draw,font=\large}
	}
}

% disponi algoritmi
\usepackage{algorithm}
\usepackage{algorithmic}
\makeatletter
\renewcommand{\ALG@name}{Algoritmo}
\makeatother

% disponi numeri di pagina
\usepackage{fancyhdr}
\fancyhf{} 
\fancyfoot[L]{\sffamily{\thepage}}

\makeatletter
\fancyhead[L]{\raisebox{1ex}[0pt][0pt]{\sffamily{\@title \ \@date}}} 
\fancyhead[R]{\raisebox{1ex}[0pt][0pt]{\sffamily{\@author}}}
\makeatother

\begin{document}
% sezione (data)
\section{Lezione del 20-03-26}

% stili pagina
\thispagestyle{empty}
\pagestyle{fancy}

% testo
\subsubsection{Esempi di variabili aleatorie discrete}
Nella scorsa lezione abbiamo visto le variabili aleatorie discrete. Vediamone adesso alcuni esempi:
\begin{itemize}
\item \textbf{\textsf{Variabili aleatorie di Bernoulli}} \\
Consideriamo un esperimento casuale che assume uno dei 2 valori possibili:
\begin{itemize}
	\item Il valore $1$ con probabilità $p$;
	\item Il valore $0$ con probabilità $1 - p$.
\end{itemize}
Una variabile aleatoria di questo tipo è detta variabile aleatoria di \textit{Bernoulli}:
$$
X \in \mathrm{Bernoulli}(p)
$$

\item \textbf{\textsf{Variabili aleatorie binomiali}} \\
Si ripeta $N$ volte un esperimento aleatorio nel quale un certo evento $A$ (evento favorevole) si può presentare con probabilità $p$ in ciascuna prova (prove indipendenti). Si può definire una variabile aleatoria $X$ il cui valore si identifica con il numero di volte in cui si verifica l’evento $A$ sul totale delle $N$ prove. Tale variabile aleatoria, detta numero di successi, è di tipo discreto e può assumere i valori $X_k = 0, 1, 2, ..., N$ con massa di probabilità:
$$
p_X(k) = P(\{ X = k \}) = \binom{N}{k} p^k (1 - p)^{N - k}, \quad 0 < p < 1, \ k = 0, 1, ..., N
$$

Una variabile aleatoria di questo tipo è detta variabile aleatoria \textit{binomiale}:
$$
X \in \mathrm{Binomiale}(p, N)
$$
che notiamo a differenza della variabile aleatoria di Bernoulli è definita da 2 parametri (probabilità e numero di prove);

\item \textbf{\textsf{Variabili aleatorie di Poisson}} \\
	Una variabile aleatoria è detta di \textit{Poisson} con parametro $\Lambda$ (con $\Lambda > 0$) se è discreta, definita per i valori interi positivi, ed ha la seguente massa di probabilità:
$$
p_X(k) = P({X = k}) = \frac{\Delta^k}{k!}e ^{-\Delta}, \quad k = 0, 1, 2, ...
$$
dove notiamo che $k$ non chiude ad $N$ ma prosegue ad $\rightarrow +\infty$.

Visto che andiamo su $k \in \mathbb{N}$, verifichiamo che è rispettata la normalizzazione della funzione di massa di probabilità:
$$
\sum_{k=0}^\infty p_X(k) = e^{-\Lambda} \sum_{k=0}^\infty \frac{\Lambda^k}{k!} = e^{-\Lambda}e^\Lambda = 1
$$
dallo sviluppo di Taylor dell'esponenziale.

Vediamo alcune proprietà:
\begin{itemize}
	\item Se $\Lambda < 1$, il picco di probabilità della funzione massa di probabilità si ha in $k^* = 0$;
	\item Se $\Lambda > 1$, il picco di probabilità della funzione massa di probabilità si ha in $k$ uguale alla parte intera di $\Lambda$, cioè:
$$
k^* = \left \lfloor{\Lambda}\right \rfloor 
$$
\end{itemize}

Come abbiamo già detto, chiamiamo questa variabile aleatoria di \textit{Poisson}:
$$
X \in \mathrm{Poisson}(\Lambda)
$$
che ha nuovamente un solo parametro ($\Lambda$).
\end{itemize}

\subsubsection{Dalla distribuzione di probabilità discreta alla funzione massa di probabilità}
Vediamo nel dettaglio come passare dalla distribuzione di probabilità discreta alla funzione massa di probabilità. Se la $F_X(x)$ (funzione massa di probabilità, 8.5) di una variabile aleatoria discreta è nota, è possibile dedurre si i valori assunti dalla variabile aleatoria, sia le corrispondenti probabilità

In particolare, come sappiamo, la distribuzione della variabile aleatoria discreata è una funzione a gradini: i gradini sono in corrispondenza degli $x_i$ (risultati) dove è definita la funzione massa di probabilità. In simboli: 
$$
F_X(x) = P(\{ X \leq x \}) = \sum_{i, \ x_i \leq x} P( \{X = x_i\} ) = \sum_{i, \ x_i \leq x} p_X (x_i) = \sum_i p_X(x_i) u (x - x_i) 
$$
dove la funzione gradino di \textit{Heaviside} $u(x)$, già vista nel campo delle trasformate di Fourier continue:
$$
u(x) = \begin{cases}
	0, \quad x < 0 \\
	1, \quad x \geq 0
\end{cases}
$$
ci permette di riportare gli $i, \ x_i \leq x$ a pedice a semplici $i$ sui risultati dell'esperimento. 

\subsection{Variabili aleatorie continue}
Veniamo quindi alle variabili aleatorie \textit{continue}:
\begin{definition}{Variabile aleatoria continua}
Una variabile aleatoria si dice \textit{continua} se assume un numero infinito non numerabile di valori (ad esempio, tutti i reali in un intervallo $[a ,b]$).
\end{definition}

Nello specifico, la continuità si definirà come:
$$
F_X(x^+) = F_X(x^-)
$$

Notiamo che se prendiamo questa definizione, la probabilità di trovare la variabile aleatoria ad un qualsiasi valore è nulla, cioè:
$$
P(\{ X = x \}) = 0, \ \forall x
$$

L'unica cosa che potremo prendere $\neq 0$ sarà allora la probabilità su un \textit{intervallo}:
$$
P(\{ x_1 \leq X \leq x_2 \}) \neq 0, \in [0, 1]
$$

Inoltre, prendere funzioni continue ci permette di sollevare le condizioni sulla strettezza delle diseguaglianze che definiscono i nostri intervalli, in quanto nulla cambia se prendiamo $<$ o $\leq$. Varrà quindi:
$$
P(\{ x_1 \leq X \leq x_2 \}) = F_X(x_2) - F_X(x_1)
$$

Nota una definizione di variabile aleatoria continua, possiamo definire la funzione \textbf{densità} di probabilità:
\begin{definition}{Funzione densità di probabilità}
Si indica come funzione densità di probabilita di una variabile aleatoria continua $X$ la funzione:
$$
f_X(x) = \frac{d F_X(x)}{dx}
$$
\end{definition}
cioè la \textit{derivata prima} della funzione distribuzione di probabilità della variabile aleatoria.

Sulla funzione densità di probabilità valgono le seguenti proprietà:
\begin{enumerate}
	\item $f_X(x) \geq 0$, visto che è la derivata prima di una funzione monotona decrescente;
	\item Dal teorema fondamentale del calcolo integrale, si ha:
$$
F_X(x) = \int_{-\infty}^x f_X(\alpha) \, d\alpha
$$
	\item Sempre riguardo all'integrazione, vale sugli intervalli (come avevamo parzialmente già detto):
$$
P(\{ x_1 \leq X \leq x_2 \}) = F_X(x_2) - F_X(x_1) = \int_{x_1}^{x_2} f_X(\alpha) \, d\alpha
$$
	\item Prendendo gli estremi a $\pm \infty$ nell'ultima proprietà, otteniamo la proprietà di normalizzazione:
$$
 F_X(\infty) - F_X(-\infty) = \int_{\infty}^{-\infty} f_X(\alpha) \, d\alpha = 1
$$
Notiamo che un corollario di ciò è che:
$$
\lim_{x\rightarrow-\infty} f_X(x) = \lim_{x\rightarrow\infty} f_X(x) = 0
$$
\end{enumerate}

\par\smallskip

\noindent
\textbf{\textsf{Intuizione dietro la densità di probabilità}}

\noindent
Prendiamo un intervallino $[x_1, x_2]$: 
$$
x_1 = x, \quad x_2 = x + dx
$$
con $dx$ infinitesimo. Dalla terza proprietà della densità di probabilità, si ha che:
$$
\int_{x}^{x + dx} f_X(x) \, dx = f_X(x) \, dx = P(x < X \leq x + dx)
$$
per cui:
$$
f_X(x) = \frac{P(x < X \leq x + dx)}{dx}
$$
Quindi, la densità probabilità rappresenta la probabilità (al variare di $x$) che la variabile aleatoria $X$ assuma valori appartenenti all’intervallo infinitesimo $[x, x+dx]$, diviso l’ampiezza infinitesima $dx$ dell’intervallo.

Questa proprietà è utile per misurare sperimentalmente la densità di probabilità di una variabile aleatoria passando alla frequenza relativa; infatti, se ripetiamo l’esperimento $N$ volte e contiamo il numero $\Delta n(x)$ di risultati per cui $x < X \leq x + \Delta x$, ottieniamo:
$$
f_X(x) \Delta x \approx P(x < x \leq x + \Delta x) \approx \frac{\Delta n(x)}{N}
$$
purché $\Delta x$ sia sufficientemente piccolo ed $N$ sufficientemente grande, quindi la densità di probabilità si può ottenere come segue:
$$
f_X(x) = \frac{\Delta n(x)}{N\Delta x}
$$

\end{document}
